\section{評価}

\subsection{評価の目的}
評価の目的は，後段検証器としてのEQおよびNOT\_EQの判定能力と，UNKNOWNおよびTIMEOUTの発生状況と原因を確認することである．TPL検出システム全体のprecisionまたはrecallの改善は，本評価の主張対象としない．

\subsection{データセットと評価方法}
評価には，等価および非等価なプログラムペアを含むEQBENCHを用いる．同一ソースを異なる環境でビルドしたペアに対するEQ判定と，非等価ペアに対するNOT\_EQ判定を確認する．EQ，NOT\_EQ，UNKNOWN，TIMEOUT，ERRORの件数，UNKNOWN理由の分布，ならびにTIMEOUT時の式サイズやループの有無などを集計する．
% TODO(draft): データセット版，対象件数，コンパイラ，最適化，ABI，制限時間，実行環境，再現手順を固定する．

\subsection{評価結果}

\subsubsection{同一ソース由来ペアに対するEQ判定結果}
% TODO(draft): 確定した実験成果物から集計表を作成する．

\subsubsection{非等価ペアに対するNOT\_EQ判定結果}
% TODO(draft): ground truthと判定結果の対応，反例の確認方法を記載する．

\subsubsection{UNKNOWNの発生状況と原因分類}
% TODO(draft): 正規のreason ID，decision effect，reason scopeなど集計軸を確定する．

\subsubsection{TIMEOUTの発生状況}
% TODO(draft): タイムアウト設定と対象式の特徴を併記する．

\subsection{考察}
正規化および型修正リトライが判定に寄与したケースを分析し，UNKNOWNの主要因を第4節で述べた実装上の制約と対応づける．また，対象関数の特徴およびタイムアウト設定とTIMEOUTの関係を分析し，解析限界を示すUNKNOWNとの違いを考察する．現段階では，ライブラリ単位のprecisionまたはrecallの改善を主張しない．
% TODO(draft): 結果確定後に，観測事実と解釈を分けて記述する．
